\cleardoublepage

\section{设计背景}

\subsection{绪论}
三维资产生成是计算机视觉与图形学领域的核心任务之一，
在游戏开发、影视制作、虚拟现实及数字孪生等领域具有广泛的应用价值。
随着深度学习技术的快速发展，基于扩散模型的生成方法凭借其强大的分布建模能力，
在图像生成领域取得了突破性进展，并逐步延伸至三维内容生成任务。

然而，三维生成任务相比二维图像生成面临更大的挑战：
三维物体的几何结构复杂、纹理信息丰富，
且从单视图图像恢复完整三维表示本身是一个高度欠定问题。
现有方法往往依赖单张图像作为输入，难以有效推断物体被遮挡部分的几何与材质，
导致生成结果在完整性与真实感方面存在明显不足。

近年来，多视角生成技术的兴起为上述问题提供了新的解决思路。
通过在三维生成过程中引入多视角图像作为条件，
可以为模型提供更丰富的外观与结构信息，
从而显著提升生成资产的几何完整性与材质保真度。
本文正是在这一背景下，探索将多视角生成先验与大规模预训练三维生成模型相结合的有效方法。

\subsection{相关工作}
早期在三维纹理生成方面的研究\cite{chen2023fantasia3d, lin2023magic3d, poole2022dreamfusion}
主要依赖于得分蒸馏采样（Score Distillation Sampling）来优化三维表示。
此后，越来越多的三维纹理生成方法\cite{feng2025romantex, huang2025mv, zhao2025hunyuan3d, lai2025hunyuan3d}
开始引入多视图生成范式：
首先借助预训练的文本到图像或文本到视频扩散模型合成多视图一致图像，
再依据各视图对应的相机位姿将纹理重投影至三维物体表面。
上述方法充分利用预训练图像或视频扩散模型中蕴含的丰富生成先验，在视觉质量上取得了显著提升。
然而，由于遮挡、覆盖不全以及视图间的不一致性，生成的纹理常常表现出纹理不完整、对齐错误或不连续等问题。

另一类研究\cite{chen2025lafite,xiang2025native}提出了原生三维纹理生成范式，以克服多视图投影方法的固有局限。
这些方法基于结构化的三维纹理表示\cite{xiang2025structured}，
将几何与纹理信息编码为三维网格的内在属性，
并通过联合训练三维纹理变分自编码器（VAE）与扩散模型，
实现几何与纹理特征的协同重建与生成。
与多视图方法相比，原生三维纹理生成在纹理一致性、视觉连贯性及几何对齐性方面均有所改善。
然而，由于此类方法对三维对象的先验知识较为有限，
在生成细粒度纹理细节以及推断遮挡区域外观方面仍存在明显不足，

\subsection{选题来源}
现有原生三维生成模型（如 Trellis.2）虽能生成高质量的三维几何与材质，
但受限于三维训练数据规模，模型对物体外观的先验知识相对匮乏，
在推断遮挡区域纹理及生成细粒度细节方面表现不足。
相比之下，以图像扩散模型为代表的二维生成模型在海量图像数据上训练，
积累了极为丰富的外观先验，对物体纹理、材质与光照具有更强的感知与推断能力。
因此，如何将图像生成模型中蕴含的丰富二维先验有效迁移至三维生成任务，
以弥补原生三维模型在外观推断方面的不足，成为本课题的核心出发点。

另一方面，一些多视角图像生成方法（如 MV-Adapter）已能够从单张图像出发，
合成视角一致、细节丰富的多视角图像序列，
为三维生成任务提供了现成可用的多视角先验。
然而，如何将这些多视角先验有效地注入现有三维生成模型，
同时避免破坏模型原有的生成能力，仍是一个尚待解决的问题。

基于上述观察，本文提出对 Trellis.2 进行参数高效微调，
以低秩自适应\cite{hu2022lora}策略改造其扩散网络中的注意力层，
使模型在保留原有三维生成能力的基础上，
获得对多视角条件输入的感知与融合能力，
从而实现从单张图像出发、多视角引导的高质量三维资产生成。

